%% LaTeX Paper Template, Flip Tanedo (flip.tanedo@ucr.edu)
%% GitHub: https://github.com/fliptanedo/paper-template-2022

% \documentclass[11pt]{article} %% Not for Lecture Notes
\documentclass[12pt, oneside]{report}    %% Has chapters

\input{FlipLectureMacros}       %% Lecture note formatting, load first
\input{FlipPreamble}			%% \usepackages, formatting
\input{FlipMacros}              %% Flip's standard macros
\input{FlipMacros_Comments}     %% Flip's macros for comments
\input{FlipMacros_Teaching}     %% Flip's macros for course notes
\input{Flip_listings}           %% Styling for code blocks
\input{FlipAdditionalHeader}    %% Modify this for each project
\input{FlipPreambleEnd}         %% packages that have to be at the end

%%%%%%%%%%%%%%%%%%%%%%%%%%%%
%% LECTURE NOTES SETTINGS %%
%%%%%%%%%%%%%%%%%%%%%%%%%%%%

% \linenumbers                  %% print line numbers (lineno package)
\graphicspath{{figures/}}       %% figure folder
\addbibresource{FlipBib.bib}    %% Define BibLaTeX source(s)

%% LEAVE THESE HERE 

\geometry{                      %% large margin for side notes
    paper=letterpaper, 
    hmargin={1cm,7.25cm},       %% 6.25cm space on right
    vmargin={2cm,2cm}, 
    marginparsep=.5cm, 
    marginparwidth=5.75cm
}

%% Def. full width; uses changepage package; 6.25cm to match hmargin difference;
\newenvironment{wide}{\begin{adjustwidth}{0cm}{-6.25cm}}{\end{adjustwidth}}


% Reset the sidenote number each section 
\let\oldsection\section
\def\section{%
  \setcounter{sidenote}{1}%
  \oldsection
}


\begin{document}

\newgeometry{margin=2cm}                   % plain geometry for frontmatter
\newcommand{\FlipTR}{UCR-TR-2024-FLIP-P165} % TR#, pdfsync may fail on 1st page
\thispagestyle{firststyle} 	               % TR#; otherwise use \thispagestyle{empty}
\pagenumbering{gobble}                     % no page number on first page 

%%%%%%%%%%%%%%%%%%%%%%%%
%%%  FRONTMATTER    %%%%
%%%%%%%%%%%%%%%%%%%%%%%%


\begin{center}
    {\large \textsf{UC Riverside Physics 165, Spring 2024} \par}
    {\huge \textbf{Introduction to Particle Physics} \par}\vspace{.5em}
    {\large {Kinematics and Dynamics of the Standard Model} \par}
    \vskip .5cm
\end{center}

\input{FlipAuthors}

\vspace{2em}\noindent
An introduction to elementary particle physics: the study of the fundamental constituents of matter and the forces that dictate their interactions. We focus on building a theoretical understanding of the Standard Model of particle physics based on Feynman diagrams. 

\vspace{5em}
\begin{center}
\includegraphics[width=.3\textwidth]{figures/P5BW.pdf}  
\end{center}  


% \vspace{2em}
\vspace*{\fill}

\noindent
\textsf{Last Compiled: \today}

\noindent
\textsf{Image: \acro{P5} Logo, Sandbox Studio}


\noindent
\textsf{CC BY-NC-SA 4.0}~\ccbyncsa 

\noindent % Course notes URL
% \url{https://github.com/fliptanedo/P231-2023-Math-Methods}

%% Front page logos
% \vspace*{\fill}
\begin{center}
\includegraphics[height=.1\textwidth]{figures/FlipAmbigram.png}
\hspace{5em}
\includegraphics[height=.1\textwidth]{figures/UCRPnA_banner.png}
\end{center}

\newpage

\small
\setcounter{tocdepth}{2}
\tableofcontents
\normalsize
\clearpage
\restoregeometry        %% Return to lecture note geometry 
\pagenumbering{arabic}  %% Turn on regular page numbers


%%%%%%%%%%%%%%%%%%%%%
%%%  THE CONTENT  %%%
%%%%%%%%%%%%%%%%%%%%%

% % \chapter{Things I'm working on}

% % \url{https://github.com/fmarotta/kaobook/issues/15}
% % % Because sidecite dosn't work
% % %\sidecite{Feng:2016ijc}

% % \input{examples_lecture}

% % % \chapter{Paper examples}

% % % Here are the standard examples I use for my \texttt{paper} template. I include them here to check that nothing has broken. These do not make use of the margin at all. You can see what happens when some text spills into the margin unintentionally.

% \input{examples}
% \input{examples_teaching}
% % % \input{examples_listings}
% \input{examples_bestpractices}
% \input{examples_refs}
% \input{examples_lecture}

% % %% CHAPTER SUBAPPENDIX %% if using report class
% % \begin{subappendices}
% % \section{Subappendix}\label{sec:subappendix:eg}
% % This chapter has its own special appendix.
% % \end{subappendices}

\chapter{The Course}

\section{Our Goal}

The goal of this course is to teach the \emph{theoretical framework} of particle physics. The underlying structure of this discipline is called \textbf{quantum field theory} and is the union of special relativity and quantum mechanics. Over one-quarter course we tell the story of leptons, quarks, and various gauge bosons---but what I \emph{really} want to convey to you is how quantum field theory works. 

Ordinarily, quantum field theory is a graduate-level course that you take after taking not only upper division quantum mechanics, but graduate-level quantum mechanics, electrodynamics, statistical mechanics/field theory, with smatterings of courses that hammer home special relativity (perhaps in general relativity course) and bits of complex analysis. Even then, the course has a reputation for being challenging because it demands a level of physical sophistication to appreciate.\sidenote[][-5em]{Usually we tell our graduate students to expect to take quantum field theory a few times in order to prepare to spend the rest of their research careers continuing to chip away at the frontier of human knowledge in this field. I am reminded of the [lightly paraprhased] quote by theorist Nima Arkani-Hamed: ``You can learn quantum mechanics in a few weeks. But you need to dedicate a lifetime of research to hope to understand quantum field theory.''} 

This course is an attempt to give a working knowledge of the big picture. You will see how Feynman diagrams are both a perturbative expansion of a transition amplitude \emph{and} a useful mnemonic for the story of particle scattering. You will see how indices are a physicist's crutch to mathematically implement symmetries. Along the way, you may come to appreciate corners of quantum mechanics and relativity that otherwise may slip the usual undergrad curriculum. 

Most of call, \emph{this course is a bridge}. For those who are interested in theoretical physics of any type\sidenote{In some sense, all theoretical physicists are quantum field theorists. Yes, even condensed matter theorists.} or those interested in particle physics of any type\sidenote{Including experimental particle physics and astro-particle physics.}: I want you to leave this course knowing how one \emph{uses} quantum field theory in particle physics so that if you go to graduate school, you will already have the bird's-eye-view of how different technical ideas come together. Alternatively, for those who know that their passions are not in this discipline: I want you to appreciate what the scaffolding of quantum field theory buys us as physicists, so that when you go off to become your future self, you are an informed ambassador of physics.

\section{What we miss}

This is \emph{not} a `modern physics course---what I mean by that are courses that are glorified ``physics for poets'' courses that offer ideas without mathematical rigor. Instead, this course is \emph{all about} understanding the mathematical rigor, even though we will not derive every step (leaning instead on analogies as appropriate) and even though the purpose of the course is not to see who can calculate the most tedious cross section.

A course like this can span multiple terms and focus on many different aspects of particle physics. Given that we have \emph{one} term and that we want to start by assuming the bare minimum, we have to make deliberate cuts to what we investigate. Rather than trying to do a little bit of everything, we will dive deeply into the theory and sacrifice the following:
\begin{itemize}
    \item The experimental foundation of particle physics. This is the biggest sacrifice because physics is ultimately an \emph{empirical} science and practitioners need to be grounded in experiment.\sidenote{Even theorists.} We excise this aspect in part because I would struggle to do it justice, but also because we are in a moment where the types of experiments that particle physicists do has evolved rapidly over the last decade.\sidenote{\tacro{UCR} students are encouraged to reach out to members of our experimental particle cosmology group to learn more about this.}
    \item The history of particle physics. I did not appreciate this as a student, but good physicists understand how those who came before them had their key \emph{aha!} moments: what what puzzles where they thinking about, how did they make progress?\sidenote{This is different from a hagiography of physics heroes. Most of our physics heroes are flawed human beings and most of the heroism is rooted in a broader collective of people than our stories usually tell.}
    \item Computational tools. Particle physicists were the original `big data' scientists with a huge throughput of data in particle colliders. Experience with some of the computational tools are a great way to get into undergraduate research in this field. I regret that this is something that we cannot fit into this course, but I encourage those interested in this field to be prepared to do computational work.
\end{itemize}

\section{Prerequisites}

I have done my best to minimize the prerequisite knowledge for this course. At the bare minimum, we require the following:
\begin{enumerate}
    \item This means you have had first-year physics and did well.\sidenote{At the very least, if you took first-year physics now you would ace it.}
    \item Analytical mechanics at the level of having some familiarity with Lagrangian mechanics and variational principles.
    \item At least two quarters of quantum mechanics so that you are familiar with bras, kets, operators, superposition, and amplitudes. You should be comfortable with angular momentum and spin.
    \item Some introduction to special relativity so that you are familiar with the principles of length contraction and time dilation.
    \item Linear algebra at the level of Physics 17 at \acro{UCR}.\sidenote{The course notes for Physics 17 may be a useful reference.\footnotemark}\footnotetext{\url{https://sites.google.com/ucr.edu/physics017/}}
\end{enumerate}

In a perfect world, you would also have the following background:
\begin{itemize}
    \item A solid background in linear algebra and some familiarity with how this relates to the representation of groups. It would help if you do not cringe at the word \emph{tensor}. It would be even better if you knew what makes a tensor a tensor as opposed to a multi-dimensional array of numbers.
    \item Some familiarity with the notion of a generator of a transformation and its exponentiation, for example from a quantum mechanics course.
    \item An idea of what a Green's function is.\sidenote{The course notes for Physics 231 may be a useful reference.\footnotemark}\footnotetext{\url{https://sites.google.com/ucr.edu/p231/}} We will not use this word too often, but understanding what it means usually correlates to a mathematics and physics background that will be helpful.
    \item Some introduction to relativistic quantum mechanics; for example, knowing the Dirac and the Klein-Gordon equations.
    \item Know what a cross section is from a mechanics course or, even better, Fermi's Golden Rule from a quantum mechanics course.
\end{itemize}
We do not live in a perfect world and it is much better to take the pioneering spirit of jumping in. As one of my undergraduate mentors told me, \emph{enthusiasm makes up for a lot of things.}\sidenote{Enthusiasm can be measured by how much time do you set aside to learn the \emph{extra} stuff that you need to best appreciate this course.}

\section{What it takes to succeed in this course}

\paragraph{Do your homework} I tell the Physics 39 classes that the best advice I can give is to \emph{do your homework}.\sidenote{Whenever I do this, I think about the Mary Schmich essay and Baz Luhrmann song, ``Wear Sunscreen.'' The song plays in the back of my head every time I talk about homework.\footnotemark}\footnotetext{\url{https://en.wikipedia.org/wiki/Wear_Sunscreen}} Trust me, I get no particular joy assigning or reviewing your homework---unlike many other topics I may teach, I actually know this shit. The value of homework is \emph{practice}.\footnote{\url{https://www.youtube.com/watch?v=p-BR1mXwtB0}} This is my commitment to you to cobble together an opportunity for \emph{you} check your understanding, to learn more (and more meaningfully), and to develop your own style of \emph{doing physics}. Part of my commitment to you is to create a space---our class---where we can work together on these meaningful exercises. I encourage you to do the homework, this is where you \emph{become a physicist}.

\paragraph{Ask questions} This course can be technical, it draws on different branches of physics, and even then it is just scratching the surface of some of the biggest open questions in science. If you engage with the material deeply enough, I \emph{expect you to be confused}. I am often confused. Do everyone a favor and ask during class. If you are embarrassed, I suggest phrasing your question as follows:
\begin{quote}
Is it obvious that...?
\end{quote}
When you ask it this way, the question is no longer about you being confused, it is about what is the most insightful way to think about a topic.

\paragraph{How to answer questions} The best part of a class is the community that we build together. I may be the person in the class who is most familiar with this material, but you and your classmates are the ones who are familiar with the journey that you're on---and there is a lot that I learn from your journey. As such, expect me to \emph{ask you questions}. This can be scary, but know that when I ask a question, I am not testing \emph{you} and what you have learned, I am testing \emph{me} and what I need to teach. If you find yourself on the business end of a question during class, consider the following options:
\begin{itemize}
    \item If you know the answer, give the answer and a justification.
    \item If you are not sure but think you know the answer, give the answer and a justification. You can even say that you are not sure and give reasons to be skeptical.\sidenote{This is part of the scientific method.} This is useful for me as well.
    \item If you have no idea what the answer is, then say so and explain what is confusing about the question. If you do not understand the question, then say that you do not understand the question and then ask a specific question for clarification. 
\end{itemize}% be ready to answer questions
% be critical: why are we doing this
% do your homework \url{https://www.youtube.com/watch?v=sTJ7AzBIJoI}

\section{These notes}

This is a first draft of these lecture notes. You can expect both outright errors and explanations that may not be fully satisfying. Please bring up any questions to me at your soonest convenience. 

I also have a tendency towards marginalia---sidenotes and footnotes. It is a habit I pick up from my favorite textbook authors (Tony Zee in particular); I feel like it brings a bit of the flavor of the course into the text. They also reflect my occasionally scatter-brained enthusiasm for this subject. If you are in a particular hurry, you can safely skip the side notes. I have tried to include references on the page that we use them rather than in a comprehensive bibliography in the back. Students who are especially dedicated to this subject are encouraged to pursue as many of these references as they are able to. 

\section{References}
Unfortunately there is not a perfect book on particle physics at this level.\sidenote{Jeff Richman at \acro{UCSB} is working on one that may be close. At the time of this writing that book is not yet complete.} Here are a few that you may consider. 

\subsection{Particle Physics}
\begin{itemize}
    \item \cite{Griffiths:2008zz}
    \item \cite{Larkoski:2019jnv}
    \item \cite{Peskin:2019iig}
    \item \cite{Cahn:1989by}
    \item \cite{Goldberg:2017dlc}
    \item \cite{Bettini:2008zz}
\end{itemize}

\subsection{Introductions to Quantum Field Theory}
Recently there has been an explosion in the number of ``quantum field theory for undergraduates'' textbooks. Here are a few that I think are particularly effective for further study. The books all start at the level assumed for this course, but each goes much deeper into the subject and wold be well suited for anyone interested in pursuing theoretical research in graduate school.
\begin{itemize}
    \item \cite{Schwichtenberg:2018dri}
    \item \cite{schwichtenberg2020no}
    \item \cite{Lancaster:2014pza}
    \item \cite{donoghue2022prelude}
    \item \cite{Zee:2010qce}
    % \item \cite{Perkins:2003pp}
    \item \cite{Feynman:1986er}
    \item \cite{Veltman:1994wz}
\end{itemize}

\chapter{Introduction}

This is a course on the theory of the Standard Model of particle physics. It is a theory so successful that it is called the \emph{Standard Model} with capital letters. Among its crowning achievements is the discovery of the Higgs boson~\autocite{CMS:2012zhx, ATLAS:2012ae} and the verification of the electron's anomalous magnetic moment~\autocite{Fan:2022eto}---the most precisely predicted and measured number in nature.\sidenote{The error is less than a part per trillion.}

\section{What is Particle Physics}

Particle physics has had a few different names. It used to be called \emph{subatomic physics} because it is the study of things that are smaller than the atom. But then nuclear physics split off as a discipline\sidenote{By the end of this course, you may come to appreciate why this is the case.\footnotemark}\footnotetext{Nuclear physics works with \emph{effective theory} of hadrons at the \GeV{} scale. This theory is distinct from the Standard Model because it is precisely the regime where the theory of quarks is non-perturbative and alternative descriptions are necessary.} and what does it mean for a particle to be `smaller' anyway?

\begin{exercise}
If someone who is not a physicist asked you what the size of the electron is, what would you say? Go ahead and look up the radius of an electron. What does this number mean? What about a photon? What is the `size' of a quantum of light?\sidenotemark
\end{exercise}\sidenotetext[][-3em]{Confused? Good. You're thinking.}

Another historical name for particle physics is \emph{high-energy physics}. This is because our experimental apparatuses are traditionally colliders where we smash together particles at high energies. 
\begin{exercise}
What does it mean for a particle to have a ``high energy''? What is this energy being compared to? The kinetic energy of a mosquito flapping its wings is negligible for most human-scale activities, but it is quite substantial when smashing together two electrons.
\end{exercise}
This is because high energy colliders are actually microscopes. Go ahead and review Rutherford scattering in your favorite quantum mechanics book. Rutherford scattering is the famous experiment that lead to the discovery of nuclei by shining an electron beam onto a gold foil target. 
\begin{exercise}
How is the energy of the electron beam related to the spatial resolution of the experiment?
\end{exercise}
In this way, one may be forgiven for conflating particle physics with \emph{collider} physics. Indeed, particle physics is still a field that looks to the Large Hadron Collider as one of its Meccas. However, unlike the case even twenty years ago, modern particle physics is far more diverse than `just' high-energy [collider] physics.\sidenote[][-5em]{Theorist Hitoshi Murayama once said that this field stretches ``\emph{from the smallest scales to the largest, from the heavens to the hell}.'' By this he was referring to the most fundamental particles, the cosmological evolution of the entire universe, space-based telescopes, and underground detectors.} As of the time of this writing, a pretty good summary of the field is the 2023 Particle Physics Project Prioritization Panel (\acro{P5}) report.\footnote{\url{https://www.usparticlephysics.org/2023-p5-report/}} The report is the guiding document for the Department of Energy and National Science Foundation---the primary funding agencies for particle physics in the \acro{US}---for the key questions that particle physics seeks to answer over the coming decade.

So I am calling this field \textbf{particle physics}. Even this name is debatable: while the objects that we \emph{observe} are \emph{particles}---usually energy and momentum eigenstates, perhaps described by some sort of creation operator $a^\dag$---the underlying mathematical objects in our theory are \emph{quantum fields}. Even then, some phenomena in particle physics seem to probe the decidedly `wavey' limit of a quantum field.\footnote{See, for example, wave dark matter, \arXiv{2101.11735}.} You could call this ``applied quantum field theory,'' but applications of quantum field theory are much broader than particle physics. So for our purposes, let's call it particle physics and not get too hung up on minutiae.

So what is particle physics? \begin{quote}In this course, particle physics is the branch of science that seeks to build a model quantum field theory that describes the elementary constituents of nature and their fundamental interactions.\end{quote} We should parse a few key words here:
\begin{itemize}
    \item Quantum field theory is a theory-of-theories. The Standard Model of particle physics is \emph{a} quantum field theory in the same way that the hydrogen atom is a model quantum mechanical system. 
    \item We are building \emph{models}. A model is a mathematical description of an actual system. The mathematical description is usually idealized and hence simplified. The model is a way to learn about how underlying principles (usually symmetries) can lead to particular phenomena. See Fig.~\ref{fig:MilleniumFalcon}.
    \item We talk about the \emph{elementary constituents} of nature. In our model, these are quantum fields. Excitations of these fields are what we call elementary particles. The name is meant to evoke some sense of indivisibility, but this is not at all necessarily the case.\sidenote[][-3em]{I once read a story about a woman on a plane who was bragging about how brilliant her son is. When she saw her seatmate reading a book, \emph{Elementary Particle Physics}, she said ``hmph! Well, my son studies \emph{advanced} particle physics.''} For example, it can be convenient to write models where protons and neutrons are elementary, even though we know that they are not. 
    \item Finally, the \emph{fundamental interactions} between particles is conventionally what are called forces. We shall see in this class in particle physics, there is no strict distinction between particle and force---they are all described by quantum fields whose excitations are particles.\sidenote{In particle physics, everything is a particle. In contrast, to astronomers, subatomic scientists are either particle physicists or a metallurgists.\footnotemark}\footnotetext{\url{https://en.wikipedia.org/wiki/Metallicity}} 
\end{itemize}
\begin{figure}%[th]
    % \centering
    \sidecaption[][-2\baselineskip]{%
        What is the difference between these two? One is a screenshot of the Millennium Falcon from \emph{Wookiepedia}, the other is a picture of a \emph{model} of the Millenium Falcon from Lego. Physicists build models of nature. We hope that the model may illuminate principles of how nature works, but we do not confuse the model as a prescription for how nature must behave.
        \label{fig:MilleniumFalcon}
    }
    \includegraphics[width=\textwidth]{figures/MillenniumFalcon.png}
\end{figure}
These ingredients come together through a tool called a Feynman diagram.

\section{Diagrammar}

Fig.~\ref{fig:ee:gamma:gamma:example} is an example of a Feynman diagram.
\begin{marginfigure}%[th]
    \includegraphics[width=.8\textwidth]{figures/feyn_eegaga.pdf}
    \captionsetup{font={scriptsize,sf}}
    \caption{Example of a Feynman diagram. Here an electron and a positron annihilate into a pair of photons.}
    \label{fig:ee:gamma:gamma:example}
\end{marginfigure}
Feynman diagrams are a perturbative expansion of the quantum mechanical amplitude for a something to happen. You may recall that the quantum amplitude is a ``square root of the probability\footnote{This notion comes from quantum mechanics where the probability of observing a state $\ket{\Psi}$ is $\langle \Psi | \Psi\rangle = |\Psi|^2$. In quantum field theory we typically talk about \emph{cross sections}, $\sigma$. The relation between the amplitude $\mathcal M$ and cross section is $\sigma \sim |\mathcal M|^2$ and contains kinematic factors. We determine these kinematic factors in subsequent chapters.}.'' The statement that diagrams are a \emph{perturbative expansion} means that there is some small parameter for which we are performing a Taylor expansion\sidenote{Formally the amplitudes are complex and may contain singularities. It is thus more appropriate to say that this is a \emph{Laurent expansion}. One of the `deep' ideas in quantum field theory is the intimate relationship of \emph{analyticity} (complex differentiability) to physical properties. To dig deeper, see my Physics 231 notes.\footnotemark}\footnotetext{\url{https://sites.google.com/ucr.edu/p231/}}. In the standard case, this perturbative expansion is usually an expansion in couplings.

By \textbf{coupling} I mean a parameter of the theory that determines how much some particles interact with each other. When the coupling is large, the interaction is very strong. When the coupling is small, the interaction is very weak. One coupling that you may be familiar with is the electrodynamic coupling, $e$. You are probably most familiar seeing $e$ as an ingredient in the fine structure constant, 
\begin{align}
    \alpha = \frac{1}{\hbar c \varepsilon_0} \frac{e^2}{4\pi} \approx \frac{1}{137} \ .
\end{align}
The first factor of $(\hbar c \varepsilon_0)^{-1}$ are relics of using silly units. When we use natural units---see Sec.~\ref{sec:units:dimensions}---these are set to one. You can see that $1/137$ is a small number, so it at least makes sense that if we had some amplitude $\mathcal M(\alpha)$ that is a function of the electrodynamic couplings through $\alpha$ that we could imagine doing the perturbative expansion
\begin{align}
    \mathcal M = \mathcal M_0 + \alpha \mathcal M_1  + \alpha^2 \mathcal M_2 + \cdots \ ,
    \label{eq:amplitude:perturbative:expansion}
\end{align}
and then dropping any subleading terms since we expect them to be percent-level corrections. If the couplings are large, then this expansion breaks down because subsequent terms are not small. In fact, this is what happens with the strong interactions (quantum chromodynamics), the force that holds nuclei together. Thus there are regimes where the usual Feynman expansion fails: it seems like we should not use these diagrams to describe the interactions of the quarks and gluons that are `inside' a proton.
\begin{example}
I seem to have implied that Feynman diagrams do not work for the strong interaction. Despite this, collider physicists working on the Large Hadron Collider `speak' the language of Feynman diagrams. They'll even draw diagrams that involve the strong force. What gives? Apparently I haven't told you the whole story...
\end{example}
Note that I wrote \emph{couplings} not the common phrase \emph{coupling constants}. That is because---brace yourselves---these couplings are generally \emph{not} constant. In fact, they depend on the energy scale at which you probe them. If I smash together color-charged particles\sidenote[][]{In this class and in this field, we write \emph{colored} to mean color-charged, or charged with respect to the strong force. See Chapter~5 of \emph{The Disordered Cosmos} for an anthropological discussion.\footnotemark}\footnotetext{\cite{prescod2021disordered}} at high energies, the analogous fine structure parameter for the strong force, $\alpha_\textnormal{s} = g_\textnormal{s}^2/4\pi$ depends on the characteristic energy scale\footnote{We write $\sqrt{Q^2}$ rather than $E$ because $Q^2$ is a Lorentz-invariant quantity, as we explain below in our review of special relativity.} $\sqrt{Q^2}$ at which one probes the interaction, see Fig.~\ref{fig:aS:running:from:1604.08082}. 
\begin{marginfigure}%[th]
    \includegraphics[width=.8\textwidth]{figures/aSrun_1604.08082.png}
    \captionsetup{font={scriptsize,sf}}
    \caption{Approximate value of the strong force fine structure parameter, $\alpha_\textnormal{s}$. Lines correspond to slightly different calculations. The horizontal axis is the square of the characteristic energy scale at which $\alpha_\textnormal{s}$ is being measured. From Fig.~3.2 of \arXiv{1604.08082}.}
    \label{fig:aS:running:from:1604.08082}
\end{marginfigure}
This idea should be shocking the first time that you see it. The `coupling constants' are not constant at all. The meaning of this oddity is an idea called \emph{renormalization} and is rooted in making sense of what the \emph{actual} small parameter is in perturbation theory: it turns out not to be the fine structure constant, but the fine structure constant times a function of the kinematics of a process. This means that at sufficiently high energies, we can meaningfully talk about Feynman diagrams of quarks exchanging gluons. However, at low energies, these diagrams lose their meaning. If you were paying attention in the previous sub-section, this is the regime where particle physics becomes nuclear physics.
\begin{example}
This reminds me of the difference between a \emph{physicist} and a \emph{physics fan}. A physics fan is someone who thinks that $1/137$ is a fundamental constant of the universe and so should be tattooed on their body. A physicist stops to think: \emph{Where does this number come from?} and proceeds to learn about the scale-dependence of the electric coupling. In fact, in this course you will find that the electric coupling is not even a fundamental parameter but a combination of more fundamental parameters. I do not care whether or not you are a physics fan, but my goal is to train you as a physicist.
\end{example}

Feynman diagrams are graphs that represent a mathematical expression for a complex number. The graphs are trajectories in spacetime that we read from left to right. Each line in a diagram represents a particle. The lines may have decorations or labels that indicate their identity. In Fig.~\ref{fig:ee:gamma:gamma:example} the two lines on the left are an electron and a positron. The two wiggly lines on the right are each photons. Each vertex (intersection of lines) represents a factor of the coupling. In Fig.~\ref{fig:ee:gamma:gamma:example} the two vertices mean that this diagram contains two powers of the electric coupling, $e$. Thus this diagram contributes to the $\mathcal O(e^2) = \mathcal O(\alpha)$ term in the expansion of the amplitude $\mathcal M(e^+e^-\to \gamma\gamma)$ in \eqref{eq:amplitude:perturbative:expansion}. The notation $\mathcal M(e^+e^-\to \gamma\gamma)$ simply means the amplitude for an electron $e^-$ and a positron $e^+$ to turn into two photons, $\gamma\gamma$.

Something else that may be familiar from quantum mechanics: amplitudes sum together. This is the origin of quantum interference and all of the fun parts of quantum physics. Our perturbative expansion \eqref{eq:amplitude:perturbative:expansion} is one example of such a sum. However there are usually multiple diagrams that contribute at a given order in perturbation theory. One of the brilliant things about these diagrams is that they have a complementary interpretation: as a \emph{sum over histories}. This is an idea that we emphasize in our review of quantum mechanics, but the idea is this: 
\begin{quote}
The amplitude to go from some initial state to some final state is represented by the sum of all Feynman diagrams that connect the initial state to the final state. Each individual Feynman diagram in the sum represents a possible history that the initial state could have taken to reach the final state.
\end{quote}
This sum over histories interpretation is outrageous the first time you see it but is ultimately an extension of the principle of least time that underlies Snell's law in optics. The generalization to a principle of least action is the core of Lagrangian mechanics and was in fact Feynman's Ph.D thesis\autocite{feynman2005feynman}. 

The diagrams are simply a tool. Feynman's long time physics rival and co-Nobel prize winner, Julian Schwinger, was a master of calculating amplitudes in quantum field theory \emph{without} any diagrams. In one interpretation of history, Feynman's main contribution at this stage of physics history was bringing that calculational technology to the masses by giving each term an intuitive meaning.\sidenote{It is no surprise that he is also (mis-)attributed as the spokesperson of the ``shut up and calculate'' school of quantum physics.\footnotemark}\footnotetext{\cite{10.1063/1.1768652}} Veltman has a nice physics-oriented history of the spread of Feynman diagrams in his book \emph{Diagrammatica}\autocite{Veltman:1994wz}. This book should not be confused with a set of lecture notes he co-wrote with 
't~Hooft called ``Diagrammar\autocite{tHooft:1973wag},'' which inspired the name of this section.
% Not for this class, but also interesting: https://arxiv.org/abs/2109.06889
% Diagrammar of physical and fake particles and spectral optical theorem
%  Anselmi


As a final note: there is far more to quantum field theory than Feynman diagrams. As a perturbative expansion, Feynman diagrams represent the \emph{easiest} part of quantum field theory. But just as there are functions that cannot be meaningfully Taylor expanded\sidenote{Try expanding $\exp(-x^{-2})$ around the point $x_0 = 0$. Every term at every order vanishes, even though the function is well defined and non-zero away from the origin. The quantum field theory analog is something called an instanton and is beyond the scope of this course.}, there are vast swaths of field theory that are intrinsically \emph{non-perturbative}. Quantum chromodynamics at low energies---where we must turn to methods in nuclear physics---is one example. We make this caveat to emphasize that although we lean heavily on the diagrammatic interpretation of particle physics, we are still just scratching the surface.



% HW: muon shot. 
% Idea: every homework is a story



\section{Kinematics and Dynamics}



\section{Units and Dimensional Analysis}
\label{sec:units:dimensions}

\section*{Acknowledgments}

\acro{PT}\ thanks 
all the people who taught him quantum field theory and particle physics over the years. In particular, courses from Scott Thomas, Pat Burchat, Savas Dimopoulos, Aaron Roodman, Michael Peskin, Shamit Kachru, David Tong, Maciej Dunajski, Ben Allanach, Hugh Osborn, Fernando Quevedo, Silvia Pascoli, Csaba Cs\'aki, Maxim Perelstein, and Yuval Grossman. I am further indebted to those who were (and are still) on this journey to figure this all out.
%
% \acro{PT} thanks 
%     the Aspen Center for Physics (\acro{NSF} grant \acro{\#1066293})
%     % and the Kavli Institute for Theoretical Physics (\acro{NSF} grant \acro{PHY-1748958})`'
%     for 
%     its 
%     % their
%     hospitality during a period where part of this work was completed. 
% %
% \acro{PT} is supported by the \acro{DOE} grant \acro{DE-SC}/0008541.
\acro{PT} is supported by a \acro{NSF CAREER} award (\#2045333).

%% Appendices
\appendix
% \chapter{Proper appendix}
% index that follows this chapter.

% \section{Things to work on}

% It may be nice to incorporate something like \texttt{classicthesis}\footnote{\url{https://www.ctan.org/tex-archive/macros/latex/contrib/classicthesis/}}


\printindex

%% Bibliography
%% USING BIBLATEX, SKIP BIBTEX
%% Use inspireHEP bibtex entries when possible
% \bibliographystyle{utcaps} 	% arXiv hyperlinks, preserves caps in title
% \bibliography{bib title without .bib}


\end{document}